\documentclass[12pt]{article}

% Packages
\usepackage{graphicx}
\usepackage{subcaption}
\usepackage{hyperref}

\usepackage[margin=1in]{geometry}

% Translations

\ifdefined\swedish
    \newcommand{\text}[2]{#2}
\else
    \newcommand{\text}[2]{#1}
\fi

\graphicspath{{./graphics/}}

\pagenumbering{gobble}
\usepackage[utf8]{inputenc}

\begin{document}

\begin{figure*}
    \begin{subfigure}[]{0.4\textwidth}
        \includegraphics[height=5cm]{photo}
    \end{subfigure}%
    ~
    \begin{subfigure}[]{0.5\textwidth}
        \part*{Joel Oskarsson}

        \begin{tabular}{l l}
            Alsättersgatan 32 & \text{+46706724739}{0706724739}\\
            584 32 Linköping & \href{mailto:joel.oskarsson@outlook.com}{joel.oskarsson@outlook.com}\\
            \href{http://joeloskarsson.github.io}{joeloskarsson.github.io} & \href{http://linkedin.com/in/joel-oskarsson}{linkedin.com/in/joel-oskarsson}
        \end{tabular}

     \end{subfigure}%
\end{figure*}

\section*{\text{Education}{Utbildning}}
\begin{itemize}
    \item \text{
            In progress: \textbf{Exchange Year, ETH Zürich}, Zürich, Switzerland\\
        }{
            Pågående: \textbf{Utbytesår, ETH Zürich}, Zürich, Schweiz\\
        }
        Sep 2018 -- Aug 2019\\
        \text{
            First year of my master as an exchange student at ETH. Courses mainly in AI and machine learning.
        }{
            Första året av masterstudier vid ETH. Kurser främst inom AI och maskininlärning.
        }

    \item \text{
            In progress: \textbf{Master programme in Computer Science and Engineering (Swedish Civilingenjörsprogram), Linköping University}, Linköping, 300 ECTS\\
        }{
            Pågående: \textbf{Civilingenjörsprogram i Datateknik, Linköpings Universitet}, Linköping, 300HP\\
        }
        Aug 2015 -- 2020\\
        \text{
            Bachelor courses in mathematics, software development and hardware. Master focused on artificial intelligence. Mean grade of 4.95.
        }{
            Kurser på kandidatnivå inom matematik, mjukvaruutveckling och hårdvara. Masterprofil inom artificiell intelligens. Medelbetyg 4.95.
        }
        \begin{itemize}
            \item \text{
                    Bachelor thesis: \href{http://urn.kb.se/resolve?urn=urn:nbn:se:liu:diva-149042}{Real-time Multiplayer Game on IoT-backend}
                }{
                    Kandidatrapport: \href{http://urn.kb.se/resolve?urn=urn:nbn:se:liu:diva-149042}{Realtidsmultiplayerspel på IoT-backend}
                }
        \end{itemize}

    \item \text{
            \textbf{Technology Program, Lindeskolan}, Upper secondary school, Lindesberg\\
        }{
            \textbf{Teknikprogrammet, Lindeskolan}, Gymnasium, Lindesberg\\
        }
        Aug 2012 -- Jun 2015\\
        \text{
            Courses in web development and network technologies.
        }{
            Kurser inom webutveckling och nätverksteknik.
        }

\end{itemize}

\section*{\text{Previous Work}{Tidigare Arbete}}
\begin{itemize}
    \item \textbf{Summer Intern, Ericsson}, Linköping\\
        Jun-Aug 2018\\
        \text{
            Summer internship at Ericsson Research. Worked with GNSS positioning.
        }{
            Sommarjobb vid Ericsson Research. Arbetade med GNSS-positionering.
        }

    \item \textbf{\text{Teaching Assistant, Linköping University}{Kursassistent, Linköpings Universitet}}, Linköping\\
        \text{Multiple periods}{Flera perioder} 2016-2018\\
        \text{
            In charge of lessons and seminars for courses in mathematics and computer science. Developed my leadership skills and personal responsibility.
        }{
            Hållit i lektioner och seminarier för kurser i matematik och datavetenskap. Har utvecklat mitt ledarskap och personliga ansvarstagande.
        }

    \item \textbf{\text{Store Clerk}{Butikspersonal}, Hemköp}, Lindesberg\\
        \text{Summers}{Somrar} 2011-2016\\
        \text{
            Extra job. Worked under a lot of individual responsibility.
        }{
            Extrajobb. Arbetade under mycket eget ansvar.
        }

\end{itemize}

\section*{\text{Experiences}{Erfarenheter}}
\begin{itemize}
    \item \text{
            Developed multiple \textbf{hobby projects} such as smaller programs and video games. Some of these, as well as some university assignments, can be found on my \href{http://github.com/joeloskarsson}{GitHub}.
        }{
            Utvecklat många \textbf{hobbyprojekt}, så som mindre program och spel. Vissa av dessa, samt en del uppgifter från universitetskurser, hittas på min \href{http://github.com/joeloskarsson}{GitHub}.
        }

    \item \text{
            Experience with \textbf{competitive programming}. I have competed in the IMPA competition with good results and also taken part in the Nordic Collegiate Programming Contest. This has given me good training in algorithm construction and testing.
        }{
            Hållit på med tävlingsprogrammering. Deltagit i tävlingen IMPA med goda resultat, samt programmerings-SM. Givit bra erfarenheter av algoritmkonstruktion och testning.
        }

    \item \text{
            Plenty of experience \textbf{working in groups}. Good team player with experience and knowledge of different methodologies for software development.
        }{
            Stor vana av \textbf{arbete i grupp}. Erfarenhet och kunskap om olika metoder för att driva projekt, speciellt inom mjukvaruutveckling.
        }

    \item \text{
            \textbf{Board member} (IT Manager) of student association FR Ryd 2017-2018. Working with many different people and under a lot of responsibility.
        }{
            \textbf{Styrelseledamot} (IT-ansvarig) för studentföreningen FR Ryd 2017-2018. Arbetat tillsammans med många olika människor under mycket ansvar.
        }

    \item \text{
            Organized multiple \textbf{LAN-parties}, both as a member of a team and in a leading role. Part of the organizing team for D-LAN 2017 (around 200 visitors). Resulted in improved people skills and teamwork.
        }{
            Arrangerat flera \textbf{LAN-event}, både som del i ett team och i en ledande roll. Var med och arrangerade D-LAN 2017 (ca. 200 besökare). Givande erfarenheter för att arbeta med många människor och i grupp.
        }

\end{itemize}

\section*{\text{Specific Knowledge}{Kunskap}}
\begin{itemize}
    \item \text{Programming Languages and Frameworks}{Programmeringsspråk och Ramverk}
        \begin{description}
            \item [\text{Knowledgeable in}{Goda kunskaper}] C++, Java, Python \text{and}{och} Web.
            \item [\text{Experience with}{Erfarenhet av}] SciPy, scikit-learn, Tensorflow, PyTorch, React, Android, Unity, Travis CI, MATLAB, SQL, \text{and}{och} C.
        \end{description}

    \item \text{
            Speak both Swedish and English fluently and communicate well in both languages. Basic understanding of German.
        }{
            Pratar både svenska och engelska flytande och kan kommunicera väl på båda språken. Grundläggande förståelse av Tyska.
        }

    \item \text{
            Accustomed to working on both Linux and Windows environments.
        }{
            Stor vana av att arbeta i både Linux- och Windows-miljöer.
        }

    \item \text{
            Swedish driver license
        }{
            Svenskt B-körkort
        }

\end{itemize}
\vfill

% Footer of page 2
\center
\begin{tabular}{l l}
    Joel Oskarsson & \href{http://joeloskarsson.github.io}{joeloskarsson.github.io}\\
    \href{mailto:joel.oskarsson@outlook.com}{joel.oskarsson@outlook.com} & \text{+46706724739}{0706724739}
\end{tabular}

\end{document}
